\documentclass{article}
\usepackage{graphicx} % Required for inserting images

\title{VirtuLearn: Designing Interactive and Educational Experiences Across Physical and Virtual Museum Spaces}
\author{Boran Gökcen}

\usepackage[
    backend=biber,
    style=ieee
]{biblatex}
\usepackage{tikz}

\addbibresource{references.bib}

\begin{document}

\maketitle

\section{Design of the physical waiting area (Draft)}
\subsection{Design Decisions}
    Interactive museum installations have been studied in public settings \cite{hornecker_learning_2006}
\subsubsection{notes from "Learning from Interactive Museum Installations  About Interaction Design for Public Settings"}

- Most museums want throughput which means installations are designed to only include enough interactions to get the point accross
(Harrison, Minneman and Balsamo) (1)
\\
- "active media" is particularly succesful (2)
\\
- Tracking visitor's interaction with the exhibits using Data analysis and qualitative analysis (interviews, observations) (3)
\\
-participant observation oriented by ethnographic approaches (Blomberg et al., 1993) (4)
\\
-The most succesful exhibits (attracted the many visitors, long interaction) were hands-on exhibits that required active
participation (5)
\\
-How to interest different age groups(6)
\\
-Different kind of museum visitors need different kind of things for a valuable experience (7)
\\
- show to others what people inside are doing with the vr-headsets as an art of performance
\\
-in a museum and public space with many competing object a low threshold for interaction is essential the first ten seconds need to provide an incentive to continue (8)
\\
-Hornecker and Stifter (2006) suggest that the popularity of the TV Newsroom installation was partly related to visitors being able to creatively appropriate the exhibit, rather than merely following predefined interactions.
\\
-This suggests that interactive stations in the physical exhibition space could benefit from allowing visitors to experiment freely and develop their own ways of interacting with the exhibits. (9)
\\
-big screens and open space allows hosting multiple people for one exhibition space (11)
\\
-the activities like recording for observers provides enterntainment and has a elements of "performance" and therefore was very popular (12)
\\
-importance of "layering" of activities Gammon(1999). Installations can alllow different levels of engagement. (providing early success experience for short engagement and simple small amounts of information and increasing complexity of information and actitvity with extendened engagement )
the first few seconds decide on whether a user continues or turns to any competing object of attention
\\
-the most engaging interactive experiences allow for productivity, creativity and/or communication, as these are basic human motivations and entail interaction by required activity of users(13)
\\
-Supporting group interaction is essential in museum exhibitions, as many visitors attend in groups and consider the shared experience an important part of their visit. Interactive installations should therefore encourage social interaction and allow multiple visitors to participate or observe together. (14)
\subsubsection{own input}
- qualitative analysis (children, adults, people who are familiar with VR, people who are no, etc...) and quantitive if possible
\\
- and I also could categorize visitors in my work
\\
- maybe ask questions of how much the visitor were engaging with each exhibit
\\
- different types of exhibits in the waiting area should attract different types of people
\\
- für das museum mitgestaltung der objekte vielleicht farbe wählen oder so oder irgendetwas was man speichern und mit nachhause nehmen kann
\\
- evt. etwas was in der gruppe genutzt werden kann vr wird an einem screen angezeigt oder so oder webxr an einem screen
\\
- vielleicht irgendeine challange einbauen ein quiz oder so fasnachtsquiz auf der handyapp?
\\
- the first few seconds decide on whether a user continues or turns to any competing object of attention
\\
- Display the VR user's perspective on an external screen so accompanying visitors can watch, discuss, and potentially participate in the experience.
\subsubsection{notes from "Onboarding and offboarding in virtual reality: A user-centred framework for audience experience across genres and spaces"}
\\
- Onboarding is the process of helping users understand how to participate in a VR experience, both technically and conceptually. Offboarding supports their return to the physical environment.
(p. 88, lines 84–91)
\\
- VR onboarding and offboarding should not be treated as activities that happen only inside the headset. The physical environment and the time before and after the experience also matter.
(p. 93, lines 305–328)
\\
- In the Limina Immersive Cinema case study, a welcoming and comfortable environment was created for first-time VR users. The waiting area included soft furnishings, subdued lighting and other features intended to make visitors feel comfortable.(p. 90, lines 179–195)
\\
- The Limina installation used fabric screens to protect visitors from being watched by passers-by while using VR.
(p. 90, lines 190–195)
\\
- Staff at Limina explained how the headset worked and how to put it on. Visitors were told how to request assistance and were allowed to put on the equipment independently unless help was needed.
(p. 90, lines 196–202)
\\
- Staff explained the safety procedures, including how they might physically approach or touch visitors while they were wearing a headset.
(p. 90, lines 196–202)
\\
- After VR, visitors were allowed to gradually readjust to their surroundings, remain seated if necessary and spend some time in the waiting area before leaving.
(p. 91, lines 207–214)
\\
- In a contrasting VR film festival experience, inadequate instructions, unsuitable seating and lack of privacy caused physical and social discomfort for the author.
(p. 91, lines 217–242; autoethnographic account)
\\
- Different visitors have different levels of VR experience, expectations and physical abilities. Some users may be concerned about falling, motion sickness or losing awareness of their surroundings.
(p. 101, lines 721–736)
\\
- VR onboarding should involve more than technical instructions. It may also include introducing the themes of the experience, setting expectations and preparing visitors for the transition into the virtual environment.
(p. 95, lines 440–448; p. 102, lines 751–762)
\\
- Exhibition staff have three main areas of responsibility: managing and explaining VR equipment, introducing visitors to the experience and ensuring their physical and psychological wellbeing.
(pp. 101–102, lines 738–772)
\\
- The physical location of a VR installation influences the visitor experience. Privacy, seating, spatial arrangement and the relationship between physical and virtual environments should be considered.
(p. 100, lines 673–688)
\\
- The onboarding and offboarding process should be designed around individual visitor needs rather than assuming that all VR users have the same knowledge, expectations or abilities.
(pp. 102–103, lines 781–806)
\\

\subsubsection{own notes}
- a way to put bags somewhere safe to use the VR when alone?
\\
- gibt es einen weg wie leute die die VR-Brille tragen abgeschirmt werden könnten von menschen die die brille nicht tragen?
\\
- Provide a short introduction to VR before visitors put on the headset
\\
- Explain the basic controller functions using a demonstration or short video
\\
- Explain how visitors can ask for help or stop the VR experience
\\
- Consider aform of privacy for people which are wearing the headset, so they don't feel exposed to everyone who's waiting
\\
- Make sure the visitors have enough space to move freely
\\
- Make the waiting area welcoming and comfortable
\\
- Include optional "offboarding" area where visitors can remove adjust and reflect on their experience
\\
- Make evaluation if the VR introduction actually helps visitors feel prepared and understand how to interact with the virtual exhibiton
\\
\subsubsection{notes from "Families' experiments and conversations at an open-ended exhibit in a science museum"}
- How can visitors learn from an interactive exhibit without being told exactly what to do?
\\
- Visitors could experiment freely instead of following predefined instructions
\\
- Families were able to make observations, discuss their findings and come up with explanations without guidance
\\
- The researchers investigated if small interventions from museum educators improves learning
\\
- Three conditions: no guidance, educators describing observations and educators giving explanations
\\
-Providing explanations helped reduce some incorrect interpretations
\\
- Open-ended activities can encourage experimentation and discussions between visitors
\\
- Different visitors have different abilities and prior knowledge which influences how they interact with an exhibit
\\
\subsubsection{own notes}
- vistors should be able to freely explore fasnacht masks instead of just reading information about them
\\
- maybe include small questions
\\
- give visitors additional information after interacting with the object
\\
- use hints instead of directly giving the answer
\\
- different visitors might need different amount of guidance
\\
- The WebAR application could combine free exploration with short educational explanations.
\newpage
\subsection{Conceptual Layout of the Physical Waiting Area}

\begin{figure}[htbp]
    \centering

    \resizebox{\textwidth}{!}{%
        \begin{tikzpicture}[
            x=1cm,
            y=1cm,
            every node/.style={font=\small},
            station/.style={
                draw=black!75,
                rounded corners=3pt,
                align=center,
                text width=3.0cm,
                minimum height=1.35cm,
                inner sep=5pt
            },
            flowarrow/.style={
                -{Stealth[length=2.5mm]},
                thick,
                dashed,
                draw=black!60
            }
        ]

            % Room boundary
            \draw[very thick] (0,0) rectangle (13,9);

            % Station 1 - Fasnacht discovery
            \node[station, fill=blue!12] (fasnacht)
            at (3.2,7.2) {
                \textbf{Station 1}\\
                Fasnacht Discovery\\
                Cultural Quiz
            };

            % Station 2 - WebAR
            \node[station, fill=green!12] (webar)
            at (9.8,7.2) {
                \textbf{Station 2}\\
                WebAR Exploration\\
                3D Objects
            };

            % Central waiting area
            \node[
                draw=black!60,
                dashed,
                rounded corners=3pt,
                align=center,
                minimum width=5cm,
                minimum height=1.2cm,
                fill=gray!8
            ] (waiting) at (6.5,4.7) {
                \textbf{Waiting Area}\\
                Free Movement
            };

            % Station 3 - VR onboarding
            \node[station, fill=orange!15] (onboarding)
            at (2.5,2.1) {
                \textbf{Station 3}\\
                VR Onboarding\\
                Equipment / Controls
            };

            % VIRTUE VR experience
            \node[station, fill=purple!12] (virtue)
            at (6.5,2.1) {
                \textbf{VIRTUE}\\
                VR Experience\\
                Protected Area
            };

            % Station 4 - Offboarding
            \node[station, fill=red!12] (offboarding)
            at (10.5,2.1) {
                \textbf{Station 4}\\
                Offboarding\\
                Reflection / Discussion
            };

            % Entrance (left)
            \draw[white, line width=5pt]
            (0,4) -- (0,5.2);
            \draw[flowarrow]
            (-1,4.6) -- (1.2,4.6);
            \node[above] at (-0.5,4.6) {Entrance};

            % Exit (right)
            \draw[white, line width=5pt]
            (13,4) -- (13,5.2);
            \draw[flowarrow]
            (11.8,4.6) -- (14,4.6);
            \node[above] at (13.5,4.6) {Exit};

            % Optional educational activities
            \draw[flowarrow]
            (waiting.north west)
            to[out=135,in=-90] (fasnacht.south);

            \draw[flowarrow]
            (waiting.north east)
            to[out=45,in=-90] (webar.south);

            % Transition to VR onboarding
            \draw[flowarrow]
            (waiting.south west)
            to[out=-130,in=90] (onboarding.north);

            % VR experience flow
            \draw[flowarrow]
            (onboarding.east) -- (virtue.west);

            \draw[flowarrow]
            (virtue.east) -- (offboarding.west);

        \end{tikzpicture}%
    }

    \caption{
        Conceptual layout of the physical waiting area,
        including educational stations, WebAR exploration,
        VR onboarding, the VIRTUE experience, and
        offboarding with reflection.
        Not to scale.
    }

    \label{fig:waiting-area-layout}
\end{figure}

\subsubsection{Station 1}
What:
- Interactive Quiz about Fasnacht characters and their historical meaning
\\
- Where visitors can guess the answers and then learn the solutions later through the WebAR
\\
- Visitors should be able to participate alone or discuss the answers in groups
\\
Why:
- We have active participation instead of passive reading (Hornecker & Stifter, 2006)
\\
- Low threshold for interaction; short activities but immediate feedback (Hornecker & Stifter, 2006)
\\
- We have something for larger groups as most museum visitors come in groups (Hornecker & Stifter, 2006)
\\
- Indpendent exploration (Franse et al., 2021)
\\
\subsubsection{Station 2: WebAR Exploration}

What:
- QR code to open the WebAR application
\\
- Select different Fasnacht objects
\\
- Rotate, zoom and examine 3D models
\\
- Here is could also be the Hologram
\\
- AR placement through camera
\\
- Additional information about cultural and historical meaning
\\
Why:
- Hands-on activities attracted visitor engagement (Hornecker & Stifter, 2006)
\\
- Visitors should be able to explore independently (Franse et al., 2021)
\\
- Optional guidance could support understanding (Franse et al., 2021)
\\
- Different levels of engagement: quick exploration or additional information. (Layering) (Hornecker & Stifter, 2006)
\\
\subsubsection{Station 3: VR Onboarding}

What:
- Short introduction to the VR headset
\\
- Explanation of the controller functions needed in VIRTUE (maybe Video)
\\
- How to put on and adjust the headset (maybe Video)
\\
- How to ask staff for help or stop the experience
\\
- Basic safety instructions
\\

Why:
- VR onboarding should begin before entering the virtual environment (Whittaker, 2025)
\\
- Technical instructions and staff assistance are important parts of onboarding (Whittaker, 2025)
\\
- Different visitors have different levels of VR experience (Whittaker, 2025)
\\
- Clear instructions and a comfortable environment can help address visitor uncertainty (Whittaker, 2025)


\subsubsection{VIRTUE: Protected VR Area}

What:
- Dedicated area for the VR exhibition
\\
- Enough space to not endanger someone
\\
- Some visual separation from waiting visitors (important)
\\
- Staff supervision and assistance
\\
- Possible storage area for personal belongings

Why:
- Not just te VR experience also physical surroundings influence the VR experience (Whittaker, 2025)
\\
- Physical comfort, privacy and staff assistance should be considered (Whittaker, 2025)
\\

\subsubsection{Station 4: VR Offboarding and Reflection}

What:
- Good to discuss the VR experience with others
\\
- Short reflection questions about what visitors saw or learned
\\
- Optional connection to related Fasnacht objects or WebAR content
\\
Why:
- Offboarding is part of the overall VR experience (Whittaker, 2025)
\\
- Visitors should have time to readjust to the physical environment (Whittaker, 2025)
\\
- Reflection and discussion are proposed educational extensions; Whittaker does not directly test their learning effectiveness
\\
    \newpage
    \printbibliography
\end{document}
